\section{引言：从理论到实践}

上一节课（Lecture 5）从宏观层面讲解了 GPU 的硬件架构和性能优化原理。本节课将从\textbf{实操}角度出发，学习如何编写、度量和优化 GPU 上的高性能代码。核心目标有三个：

\begin{enumerate}
    \item \textbf{掌握基准测试（Benchmarking）与性能分析（Profiling）}——理解代码在 GPU 上实际的执行方式
    \item \textbf{理解内核融合（Kernel Fusion）的重要性}——用 GeLU 和 Softmax 作为案例
    \item \textbf{掌握五种实现方式的权衡}——手写 PyTorch、原生 PyTorch、CUDA C++、Triton、torch.compile
\end{enumerate}

\begin{importantbox}{本节课的核心信息}
编程模型（PyTorch、Triton、PTX）与硬件之间存在巨大的抽象鸿沟，这会导致许多``性能之谜''。基准测试帮助我们理解扩展行为，性能分析帮助我们理解 PyTorch 函数的内部机制（最终归结为 CUDA 内核调用），查看 PTX 汇编帮助我们理解 CUDA 内核的内部工作。
\end{importantbox}

课程提供了五种实现同一函数的方法：
\begin{itemize}
    \item \textbf{Manual}：用基本 PyTorch 操作手写（最慢，因为没有融合）
    \item \textbf{PyTorch native}：调用 PyTorch 内置函数（已优化）
    \item \textbf{CUDA C++}：直接用 CUDA 编写内核
    \item \textbf{Triton}：用 Python 编写 GPU 内核
    \item \textbf{torch.compile}：让编译器自动优化手写代码
\end{itemize}

\begin{knowledgebox}{与上一节课的关系}
Lecture 5 建立了 GPU 硬件的心智模型（SM、内存层次、Tiling、Fusion 等）。本节课用代码将这些概念``落地''——你将亲手编写内核，用分析工具验证理论预测，并直接观察不同实现的性能差异。
\end{knowledgebox}

\subsection{本章小结}

本课以 MLP（线性层 + GeLU 激活堆叠）作为贯穿全课的实验对象，GeLU（逐元素操作）和 Softmax（行归约操作）作为内核编写的核心案例。掌握基准测试和性能分析是编写高性能代码的\textbf{第一步}，也是\textbf{最重要的一步}。

%% ========== Section 2: GPU 复习 ==========

